La especificación detallada con todos los \textit{endpoints}, sus
parámetros, esquemas de petición y respuesta, códigos de error y
ejemplos se ofrece por tres vías complementarias:

\begin{itemize}

  \item \textbf{Fichero \texttt{SCDEE API.yaml}} adjunto al material de
  entrega del \acs{tfg}, en formato OpenAPI 3.0. Cualquier herramienta
  compatible (Postman, Insomnia, \textit{VS Code REST Client}) puede
  importarlo.

  \item \textbf{Swagger UI}, accesible durante la ejecución del sistema
  en \texttt{http://localhost:8000/api/docs/}, permite explorar la
  \acs{api} de forma interactiva y ejecutar peticiones de prueba contra
  un \dfn{backend} en marcha.

  \item \textbf{Regeneración} a partir del código fuente mediante el
  comando \texttt{python manage.py spectacular -{}-file schema.yaml}
  ejecutado dentro del contenedor de la \acs{api}. El esquema se
  mantiene siempre sincronizado con la implementación.

  \item \textbf{Acceso directo al esquema} en formato \acs{json} o
  \acs{yaml} a través de \texttt{http://localhost:8000/api/schema/}
  mientras el sistema está en ejecución.

\end{itemize}
